\documentclass[11pt]{article}
\newcommand{\LabNumber}{5}
\newcommand{\LabTitle}{Stacks}
\newcommand{\LabTopic}{Stack ADT, array-based implementation, applications}
\input{common.tex}

\begin{document}
\labheader

\section*{Objectives}
\begin{itemize}[leftmargin=*,nosep]
  \item Define the Stack ADT and its \texttt{push}, \texttt{pop}, \texttt{peek}, \texttt{size}, \texttt{isEmpty} operations.
  \item Implement a stack backed by a fixed-capacity array.
  \item Handle stack overflow and underflow with appropriate exceptions.
  \item Apply stacks to: bracket matching, expression evaluation, and reversing a sequence.
\end{itemize}

\begin{summary}[Quick Reference]
\textbf{Stack interface.}
\begin{lstlisting}
public interface Stack<E> {
    int size();
    boolean isEmpty();
    void push(E element) throws IllegalStateException;
    E pop()  throws java.util.EmptyStackException;
    E peek() throws java.util.EmptyStackException;
}
\end{lstlisting}

\textbf{Array-based push/pop.}
\begin{lstlisting}
// push: data[top++] = e;   (check capacity first)
// pop:  return (E) data[--top];
// peek: return (E) data[top-1];
\end{lstlisting}

\textbf{LIFO rule.} Last element pushed is the first one popped.\\
\textbf{Complexity.}\quad All Stack operations: $O(1)$.
\end{summary}

\subsection*{Quick Experiments (Try It)}

\begin{tryit}{Overflow and underflow}
Create an \texttt{ArrayStack<Integer>} of capacity 3. Push four elements.
What happens? Now pop from an empty stack. What exception is thrown?
Write a try-catch block that handles both gracefully.
\end{tryit}

\begin{tryit}{The JVM method stack}
Write a method that calls itself unconditionally (no base case).
Run it and observe the exception.
How does this demonstrate that the JVM uses a stack for method calls?
\end{tryit}

\section*{In-Lab Exercises (target: 2 hours)}

\begin{labexercise}{Array-Based Stack \hfill {\small \itshape (35 min)}}
Implement \texttt{ArrayStack<E>} backed by an \texttt{Object[]} of configurable capacity.
\begin{itemize}[leftmargin=*,nosep]
  \item \texttt{push} throws \texttt{IllegalStateException("Stack is full")} when at capacity.
  \item \texttt{pop} and \texttt{peek} throw \texttt{EmptyStackException} when empty.
  \item \texttt{toString()} prints bottom to top separated by \texttt{|}, e.g.\ \texttt{[1|2|3]} where 3 is on top.
\end{itemize}
Test with 6 pushes and 4 pops. Verify \texttt{size()} after each operation.
\end{labexercise}

\begin{labexercise}{Bracket Matching \hfill {\small \itshape (30 min)}}
Write \texttt{boolean isBalanced(String expr)} that returns \texttt{true} if every opening
bracket is matched by the correct closing bracket.
Support three pairs: \texttt{()}, \texttt{[]}, \texttt{\{\}}.
\begin{itemize}[leftmargin=*,nosep]
  \item Push every opening bracket; on a closing bracket, pop and check for a match.
  \item Return \texttt{false} on mismatch or if the stack is non-empty at the end.
\end{itemize}
Test with: \texttt{"([])\{)"}, \texttt{"\{[()]\}"}, \texttt{"((("}, \texttt{""}.
\end{labexercise}

\begin{labexercise}{Reverse a String \hfill {\small \itshape (15 min)}}
Using your stack, write \texttt{String reverse(String s)} that pushes each character
then pops them all into a \texttt{StringBuilder}.
Verify \texttt{reverse("hello") = "olleh"}.
\end{labexercise}

\begin{labexercise}{Evaluating Postfix (RPN) Expressions \hfill {\small \itshape (40 min)}}
Write \texttt{int evalPostfix(String expr)} where \texttt{expr} is space-separated with
single-digit integers and operators \texttt{+}, \texttt{-}, \texttt{*}, \texttt{/}.
Push digits; on an operator, pop two operands, apply, push result.
Test: \texttt{"3 4 +"} ($= 7$), \texttt{"5 1 2 + 4 * + 3 -"} ($= 14$).
\end{labexercise}

\begin{aicollab}
\textbf{Suggested prompts:}
\begin{itemize}[leftmargin=*,nosep]
  \item ``Trace the bracket-matching algorithm step by step on \texttt{`\{[()]\}'}.''
  \item ``Explain why a stack is the natural data structure for evaluating postfix expressions.''
\end{itemize}
\medskip\textbf{Verify:} Ask the AI for a postfix expression that evaluates to 42.
Trace it manually through your algorithm before running the code.
\end{aicollab}

\takehomebanner

\begin{trace}[Stack trace]
Trace the stack contents (bottom to top) after each operation:
\begin{lstlisting}
push(5); push(3); push(8); pop(); push(2); peek(); pop(); pop();
\end{lstlisting}
What is returned by each \texttt{pop()} call?
\end{trace}

\begin{trace}[Bracket matching trace]
Trace \texttt{isBalanced("(\{[\}]\})")} step by step.
For each character show: the character, the action (push/pop/mismatch), and the stack contents.
Does the method return \texttt{true} or \texttt{false}?
\end{trace}

\begin{bug}[Broken pop]
Find and fix all errors (assume \texttt{top} starts at 0 and \texttt{push} does \texttt{data[top++] = e}):
\begin{lstlisting}
public E pop() {
    if (isEmpty()) throw new java.util.EmptyStackException();
    E val = (E) data[top];
    top--;
    return val;
}
\end{lstlisting}
\end{bug}

\begin{writecode}[Infix to postfix]
Write pseudocode (or Java) for converting a simple infix expression
(single-digit operands, operators \texttt{+} and \texttt{*} only, no parentheses)
to postfix using a stack.
State the rule for when to pop vs.\ push an operator.
\end{writecode}

\vspace{4pt}
\noindent\textbf{Submission.} Save files as \texttt{lab5\_ex*.java} and submit on the course site.

\end{document}
